\documentclass{beamer}
\usepackage{ctex}
\usepackage{lipsum}
\begin{document}
\title{讲解}
\begin{frame}
    \titlepage
\end{frame}
\begin{frame}
    \tableofcontents
\end{frame}
\begin{frame}
    \frametitle{第一章}
    \framesubtitle{数学基础}
    \section{第一部分}
    \subsection{第一小节}
    \lipsum[2]
    \alert{重点标注}
\end{frame}
\begin{frame}
    \frametitle{第二章}
    \section{第二部分}
    \subsection{第二小节}
    \lipsum[2]
\end{frame}
\begin{frame}
    \frametitle{What Are Prime Numbers?}
    \begin{definition}
        A \alert{prime number} is a number that has exactly two divisors.
    \end{definition}
    \begin{example}
        \begin{itemize}
            \item 2 is prime (two divisors: 1 and 2).
            \pause
            \item 3 is prime (two divisors: 1 and 3).
            \pause
            \item 4 is not prime (\alert{three} divisors: 1, 2, and 4).
        \end{itemize}
    \end{example}
\end{frame}
\begin{frame}
    \frametitle{There Is No Largest Prime Number}
    \framesubtitle{The proof uses \textit{reductio ad absurdum}.}
    \begin{theorem}
        There is no largest prime number.
    \end{theorem}
    \begin{proof}
        \begin{enumerate}
            \item<1-> Suppose $p$ were the largest prime number.
            \item<2-> Let $q$ be the product of the first $p$ numbers.
            \item<3-> Then $q + 1$ is not divisible by any of them.
            \item<1-> But $q + 1$ is greater than $1$, thus divisible by some prime 
            number not in the first $p$ numbers.\qedhere
        \end{enumerate}
    \end{proof}
    \uncover<4->{The proof used \textit{reductio ad absurdum}.}
\end{frame}
\end{document}